% Variante : option draft — les notes de travail deviennent visibles.
% Sans draft, \notework / \noteinmargin / \worktodo / worknotes ne composent
% rien : le même corps passe de la version relue à la version de travail sans
% retouche. Le corps partagé body.tex est compilé tel quel, puis une section
% exerce les quatre formes de notes.
\documentclass[11pt]{ocots-book}
\usepackage[lang=fr, theme=ocots, draft, solutions=inline, math=analysis]{ocots}
\lstset{language=[LaTeX]TeX}

\title{Variante : notes de travail (\texttt{draft})}
\author{Prénom \textsc{Nom}}
\date{\today}

\newcommand{\ocotsvariantnote}[1]{\begin{quote}\itshape Variante : #1\end{quote}}

\begin{document}

\begin{lstlisting}
\documentclass[11pt]{ocots-book}
\usepackage[lang=fr, theme=ocots, draft, ...]{ocots}
\end{lstlisting}
\ocotsvariantnote{\texttt{draft} (absent par défaut) : les macros
\texttt{\textbackslash notework}, \texttt{\textbackslash noteinmargin},
\texttt{\textbackslash worktodo} et l'environnement \texttt{worknotes}
deviennent visibles — hors \texttt{draft} ils ne composent rien.}

\ocotscollectsolutions
% Corps partagé par toutes les variantes : aucune option, aucun préambule.
%
% C'est la démonstration de l'étape 2 : le même contenu, sans une ligne de
% différence, se compose en français ou en anglais, avec ou sans corrigés, en
% couleur ou en noir et blanc. Seule la ligne \usepackage change.

\chapter{Chapitre de démonstration}
\begin{chapterintro}
    Une introduction de chapitre, avec \ie un exemple de macro localisée et
    \enquote{un texte cité}.
\end{chapterintro}
\section{Une section}

\begin{theorem}{Un théorème}{demo}
    Énoncé.
\end{theorem}

\begin{proof}
    Démonstration.
\end{proof}

\begin{definition}{Une définition}{}
    Énoncé.
\end{definition}

\begin{remark}
    Une remarque.
\end{remark}

\begin{assumption}
    Une hypothèse, étiquetée automatiquement.
\end{assumption}

\begin{exercise}[points=4]
    Énoncé de l'exercice.
\solution
    Corrigé.
\end{exercise}

\ocotsprintsolutions

\section{Notes de travail}

Une phrase ordinaire, suivie d'une note au fil du texte \notework{vérifier
cette hypothèse avant diffusion} et d'une note en
marge\noteinmargin{à relire}.

\worktodo{compléter la bibliographie du chapitre}

\begin{worknotes}
    Tout ce bloc est exclu hors \texttt{draft} : liste de points à reprendre,
    restes de calcul, questions en suspens.
\end{worknotes}

\end{document}
